% from ega4-18-fr.tex
\oldpage[IV-4]{162}

équivaut à dire que $Y'$ est étale sur $Y$). Lorsque $Y=\Spec(A)$, où $A$ est un anneau intègre et intégralement clos, $K$ son corps des fractions, on dira encore (par abus de langage, et lorsqu'une confusion avec la terminologie de (17.3.2, (ii)) n'est pas à craindre) que $L$ est non ramifiée sur $A$ au lieu de « non ramifiée sur $Y$ ».
\end{env}

\begin{remark}[18.10.11]
\EGAReaderTarget{ega4.c04.rem.18-10-11}{IV.18.10.11}
Au lieu de dire que $L$ est non ramifiée sur $Y$ certains auteurs disent aussi que $L$ est « non ramifiée sur $K$ »; nous nous garderons de suivre cet usage qui prête à confusions.
\end{remark}

On peut alors exprimer (18.10.9) sous la forme suivante:

\begin{corollary}[18.10.12]
\label{IV.18.10.12}
Soient $Y$ un préschéma normal et intègre, $K=R(Y)$ son corps de fonctions rationnelles. Alors le foncteur $X\rightsquigarrow R(X)$ est une équivalence de la catégorie des revêtements étales de $Y$, et de la catégorie des $K$-algèbres finies étales qui sont non ramifiées sur $Y$. On obtient un foncteur quasi-inverse en faisant correspondre à toute $K$-algèbre finie étale $L$, non ramifiée sur $Y$, la fermeture intégrale $Y'$ de $Y$ dans $L$.
\end{corollary}

\begin{proposition}[18.10.13]
\label{IV.18.10.13}
Soient $Y$ un préschéma normal et intègre, $K=R(Y)$ son corps de fonctions rationnelles.
\begin{enumerate}
  \item[(i)] Le corps $K$ est une $K$-algèbre non ramifiée sur $Y$.
  \item[(ii)] Soit $L$ une extension finie de $K$, non ramifiée sur $Y$, et soit $Z$ la fermeture intégrale de $Y$ dans $L$. Alors, si $M$ est une $L$-algèbre non ramifiée sur $Z$, $M$ est aussi une $K$-algèbre non ramifiée sur $Y$ (« transitivité » de la non-ramification).
  \item[(iii)] Soient $Y'$ un préschéma normal et intègre, $K'$ son corps de fonctions rationnelles, $g:Y'\to Y$ un morphisme dominant. Si $L$ est une $K$-algèbre non ramifiée sur $Y$, alors $L\otimes_K K'$ est une $K'$-algèbre non ramifiée sur $Y'$ (propriété de « translation »). Si de plus $Y=\Spec(A)$ et $Y'=\Spec(A')$ sont affines, et si $C$ est la fermeture intégrale de $A$ dans $L$, alors $C'=C\otimes_A A'$ est la fermeture intégrale de $A'$ dans $L\otimes_K K'$.
\end{enumerate}
\end{proposition}

L'assertion (i) est immédiate, $Y$ étant par hypothèse son propre normalisé dans $K$. Pour prouver (ii), soit $Z'$ la fermeture intégrale de $Z$ dans $M$, qui est par hypothèse un revêtement étale de $Z$. Comme $Z$ est étale et séparé sur $Y$, il en est de même de $Z'$ (17.3.3) qui est évidemment la fermeture intégrale de $Y$ dans $M$. Comme $M$ est une $K$-algèbre finie et séparable, il résulte de (18.10.10) que $M$ est non ramifiée sur $Y$. Enfin, pour démontrer (iii), notons que si $Z$ est la fermeture intégrale de $Y$ dans $L$, $Z\times_Y Y'$ est étale et séparé sur $Y'$ (17.3.3), fini sur $Y'$ puisque $Z$ est fini sur $Y$, et admet pour anneau de fonctions rationnelles $L\otimes_K K'$ (I, 3.4.9); comme $Y'$ est normal, il en est de même de $Z\times_Y Y'$ (17.5.7), qui est donc la fermeture intégrale de $Y'$ dans $L\otimes_K K'$, ce qui achève la démonstration.

\begin{corollary}[18.10.14]
\label{IV.18.10.14}
Soient $Y$ et $Y'$ deux préschémas normaux et intègres, $K$, $K'$ leurs corps de fonctions rationnelles respectifs, $g:Y'\to Y$ un morphisme dominant, de sorte que $K'$ est une extension de $K$. Soient $L$ une extension finie de $K$ non ramifiée sur $Y$, $L_1$ une extension composée (Bourbaki, \emph{Alg.}, chap.~VIII, \textsection~8, déf.~1) de $L$ et de $K'$. Alors $L_1$ est une extension de $K'$ non ramifiée sur $Y'$. Si de plus $Y=\Spec(A)$, $Y'=\Spec(A')$ sont affines, et si $C$ est la fermeture intégrale de $A$ dans $L$, alors le sous-anneau $C_1=A[C,A']$ de $L_1$ est la fermeture intégrale de $A'$ dans $L_1$.
\end{corollary}

En effet, $L\otimes_K K'$ est composée directe d'extensions finies séparables de $K'$, et $L_1$

